\chapter{Metodología}\label{chap:metodologia}
La metodología se concibe para sostener reproducibilidad, interpretabilidad y validez estadística en un escenario de muestra limitada, articulando la física de adquisición con un andamiaje de inferencia probabilística y de control de incertidumbre. El flujo integra gobernanza y curaduría de datos, normalización geométrica–radiométrica con justificación óptica, extracción de descriptores con fundamento físico–estadístico, ajuste de clasificadores lineales y no lineales con regularización explícita, validación cruzada anidada y verificación de calibración, además de un subsistema de detección de entradas fuera de dominio y un régimen de documentación que preserve la trazabilidad de cada decisión técnica y de cada artefacto generado \cite{Tsamardinos2022,Vabalas2019,Pedregosa2011,EAFIT2018}.

\section{Datos, gobernanza y normalización}
El estudio emplea 304 hologramas rotulados capturados con un único montaje DLHM y almacenados en formato PNG de 8 bits, con resolución máxima de \(4640\times3506\) píxeles. La clase saludable comprende 151 interferogramas de eritrocitos sin rasgos patológicos y la clase patológica 153 interferogramas correspondientes a anemia falciforme; a efectos de evaluación de robustez fuera de dominio se dispone, en repositorio separado, de 30 hologramas de organismos no sanguíneos—esporas de \emph{Aspergillus}, diatomeas y partículas de gramíneas—que no participan del ajuste del clasificador. La gobernanza de datos asegura integridad mediante sumas criptográficas y verificación de cabeceras, mantiene un registro de control de archivos y preserva únicamente metadatos técnicos mínimos (fecha y operador) necesarios para trazabilidad y auditoría, eliminando identificadores sensibles conforme al reglamento de propiedad intelectual de la Universidad EAFIT \cite{EAFIT2018}. La organización del repositorio impone versionado de directorios y artefactos derivados, de modo que cualquier figura o métrica reportada pueda rastrearse al estado específico del conjunto y al commit de código que la produjo, alineando prácticas con recomendaciones contemporáneas de ciencia reproducible \cite{Tsamardinos2022}.

La normalización se diseña como un proceso determinista que preserva relaciones espectrales y evita ambigüedades de muestreo. Cada holograma \(I\) se corrige radiométricamente (cuando hay marcos \emph{dark/flat}) y se lleva a rango unitario mediante percentiles robustos \(P_{1}\)–\(P_{99}\), estabilizando el contraste frente a atípicos. El tamaño objetivo de \(512\times512\) píxeles se obtiene por recorte central cuando la resolución lo permite para no diluir contenidos de alta frecuencia y por interpolación bilineal en caso de resoluciones inferiores, fijando una malla de muestreo coherente con las restricciones de Nyquist tratadas en el marco óptico \cite{Greenbaum2014}. Se aplica CLAHE para realce de contraste local y un filtrado de mediana suave para atenuar ruido impulsivo sin erosionar bordes interferométricos \cite{Singh2017}. Las desalineaciones por inclinación o curvatura del frente se modelan como fases cuadráticas y se compensan en el dominio frecuencial retirando términos de baja orden, lo que mejora la estacionariedad de las franjas y reduce fugas espectrales en operadores espectrales posteriores \cite{Amann2019,OConnor2020}. Cada corrida escribe una bitácora con parámetros aplicados, hashes de entrada y salida, estadísticas radiométricas pre y post normalización, y marcas temporales; esta documentación permite reproducir fielmente el estado del sistema y detectar divergencias al incorporar nuevos lotes de datos \cite{Tsamardinos2022,Vabalas2019}.

\section{Extracción de descriptores físicos y estadísticos}
La representación de cada holograma normalizado \(I_{n}\) se construye sin reconstruir el campo complejo, preservando la física de la formación del interferograma mediante descriptores complementarios en dominios espacial, frecuencial y multiescala. La microtextura se codifica con \(\operatorname{LBP}_{P,R}\) en variantes uniformes e invariantes a rotación, cuyo histograma normalizado conserva diferencias locales de intensidad con estabilidad frente a cambios suaves de iluminación y a pequeñas deformaciones geométricas \cite{Dheepak2023}. La textura de segundo orden se resume con matrices de coocurrencia \(P_{\boldsymbol{\Delta},\theta}\) y con métricas derivadas—contraste, homogeneidad, energía, correlación—promediadas sobre orientaciones canónicas para amortiguar direccionalidades espurias, lo que captura la granularidad interferométrica y la repetitividad de las franjas \cite{Barburiceanu2021}. En el dominio de Fourier se calcula la energía por anillos concéntricos y un índice de anisotropía basado en el segundo momento angular del espectro, sensibles a la concentración de energía en bajas frecuencias y a la direccionalidad del patrón, respectivamente \cite{Amann2019,OConnor2020}. La respuesta orientada se enriquece con bancos de Gabor parametrizados por orientación y longitud de onda, cuyos momentos sobre \(|I_{n}*g_{\theta,\lambda}|\) discriminan periodicidades y oblicuidades ligadas a deformaciones típicas \cite{Alzubaidi2020}. La multiescala se incorpora mediante energías de detalles wavelet \(W^{(\ell)}\) en familias Daubechies de soporte compacto, proporcionando sensibilidad a irregularidades finas en varias resoluciones \cite{Alzubaidi2020}. Para completar con morfometría, una máscara binaria obtenida por umbralización adaptativa y limpieza morfológica habilita métricas geométricas—área, perímetro, circularidad, solidez y excentricidad—junto con momentos globales de intensidad y entropía de Shannon, de utilidad cuando cambios de fase inducen contornos efectivos persistentes en intensidad \cite{Singh2017,Alzubaidi2020}.

Todos los descriptores se estandarizan con mediana y MAD para reducir sensibilidad a colas pesadas y se filtran por varianza casi nula antes del ajuste. La reducción supervisada se efectúa con \(\mathrm{SelectKBest}\) dentro de los pliegues de entrenamiento, utilizando estadísticos acordes a la escala de cada característica; esta decisión limita dimensionalidad efectiva, controla colinealidades residuales y evita fuga de información al mantener los pasos de selección estrictamente confinados a los conjuntos de entrenamiento en el marco de validación anidada \cite{Tsamardinos2022,Vabalas2019}. La configuración de radios, distancias y orientaciones se fijó tras exploración documentada que comparó alternativas con base en separabilidad y estabilidad entre ejecuciones, manteniendo un registro de parámetros que facilita auditoría y reproducción. Los vectores \(\mathbf{x}\) y metadatos asociados se serializan con huellas criptográficas y referencias a la versión de código, asegurando consistencia entre estado de datos, configuración y resultados.

\section{Modelado, validación y calibración}
El aprendizaje se formula como minimización de riesgo empírico regularizado, con un portafolio de clasificadores que aportan sesgos inductivos complementarios. La regresión logística penalizada modela \(p(y=1\mid\mathbf{x})\) mediante un mapeo sigmoidal de combinación lineal de descriptores y emplea penalizaciones \(L_{1}\), \(L_{2}\) o elásticas para promover parsimonia y estabilidad numérica; los coeficientes estimados se interpretan en términos de log-odds y contribuyen a la trazabilidad clínica \cite{Pedregosa2011}. La máquina de soporte vectorial con núcleo RBF despliega una frontera no lineal mediante proyección implícita a un RKHS, controlando el equilibrio entre error empírico y complejidad a través de \(C\) y \(\gamma\); las puntuaciones de margen se transforman luego a probabilidades mediante calibración posterior independiente del ajuste de hiperparámetros \cite{Pedregosa2011}. Los modelos basados en árboles abordan interacciones y umbrales sin supuestos de linealidad: el bosque aleatorio reduce varianza por agregación sobre remuestreos y subespacios de características, mientras que la potenciación de gradiente construye una aproximación aditiva con árboles poco profundos que optimiza la pérdida logística con regularización funcional implícita mediante tasa de aprendizaje y profundidad \cite{Friedman2001}. Estas tres familias capturan de forma conjunta relaciones lineales interpretables, separaciones de margen en espacios no lineales y particiones jerárquicas con interacción de variables.

Las salidas probabilísticas de los modelos se integran en una mezcla convexa \(p_{\mathrm{ens}}(\mathbf{x})=\sum_{k}w_{k}p_{k}(\mathbf{x})\), con pesos aprendidos en un pliegue de calibración dedicado para minimizar la pérdida de Brier; esta combinación atenúa la varianza de estimación, rescata contribuciones complementarias y evita que un único inductivo domine en regiones de escasez de datos \cite{Pedregosa2011}. La selección de hiperparámetros—\(\lambda,\alpha\) en logística; \(C,\gamma\) en SVM; número de árboles, profundidad y tasa de aprendizaje en ensambles—se realiza mediante búsqueda aleatoria acotada dentro de una validación cruzada interna estratificada; el número \(k\) de características retenidas se optimiza de forma conjunta. La capa externa de evaluación reserva un subconjunto de prueba estratificado y reporta métricas umbral-independientes (ROC-AUC y, cuando procede, PR-AUC) y umbral-dependientes (precisión balanceada, sensibilidad y especificidad), con intervalos de confianza al 95\,\% obtenidos por \emph{bootstrap} estratificado de alta repetición; la estricta anidación garantiza que toda transformación y selección se ajuste exclusivamente con información del entrenamiento \cite{Vabalas2019,Tsamardinos2022}. La calibración probabilística se verifica con pérdida de Brier y diagramas de confiabilidad; cuando se detecta descalibración se aplica ajuste isotónico o mapeo logístico sobre un conjunto de calibración distinto del empleado para la búsqueda de hiperparámetros, manteniendo independencia entre etapas \cite{Pedregosa2011}. La decisión clínica se establece fijando un umbral \(\tau\) costo–sensible que garantice sensibilidad mínima cuando el falso negativo tiene costo elevado, o bien a través del índice de Youden cuando se busca equilibrio, documentando la incidencia de \(\tau\) en los intercambios sensibilidad–especificidad.

\section{Incertidumbre, anomalías y reproducibilidad operativa}
La comunicación de incertidumbre se articula con predicción conforme en esquema \emph{split}, utilizando como medida de no conformidad funciones de las probabilidades calibradas del ensamble. El cuantil \(q_{1-\epsilon}\) estimado en el conjunto de calibración define conjuntos de predicción \(\Gamma_{\epsilon}(\mathbf{x})\) con garantía de cobertura marginal \(1-\epsilon\) bajo intercambiabilidad, habilitando abstenciones informadas cuando la evidencia es insuficiente y alineando el sistema con prácticas clínicas de revisión prioritaria \cite{Montero2024}. La detección de entradas fuera de dominio se implementa con distancia de Mahalanobis \(D_{M}^{2}=(\mathbf{x}-\boldsymbol{\mu})^{\top}\boldsymbol{\Sigma}^{-1}(\mathbf{x}-\boldsymbol{\mu})\) calculada sobre el subespacio de características seleccionadas; la media y la covarianza se estiman mediante procedimientos robustos tipo \emph{Minimum Covariance Determinant} para resistir contaminación y colinealidades, y los umbrales \(\tau_{\alpha}\) se fijan frente a la referencia \(\chi^{2}_{d}\) con validación por \emph{bootstrap} \cite{Mahalanobis1936,Rousseeuw2019}. La descomposición por coordenadas de \(D_{M}^{2}\) aporta contribuciones específicas de cada descriptor, lo que facilita la interpretación de alertas y el diálogo con especialistas. El conjunto externo de 30 hologramas no sanguíneos se reserva para valorar tasas de verdadera y falsa alarma del detector y para verificar estabilidad de umbrales frente a variaciones de dominio. 

La reproducibilidad operativa descansa en una automatización integral del flujo a través de una interfaz de línea de comandos que parametriza rutas de datos, configuración de extracción, selección de características, esquemas de validación y calibración. Cada ejecución genera artefactos versionados—vectores de características con metadatos, modelos y calibradores serializados—además de reportes con curvas ROC/PR, matrices de confusión, diagramas de confiabilidad y análisis de estabilidad de selección, todos con huellas criptográficas y referencias explícitas a la versión del código \cite{Tsamardinos2022,EAFIT2018}. La coordinación periódica con el equipo clínico permite revisar casos residuales, recalibrar umbrales costo–sensitivos y ajustar el balance entre sensibilidad y especificidad; la fase piloto prevista se ejecutará en paralelo con el flujo diagnóstico estándar para recolectar evidencia prospectiva, actualizar calibraciones y analizar desempeño por subgrupos antes de evaluar cualquier despliegue asistencial \cite{Montero2024}. En este marco, la modularidad del código habilita la incorporación futura de descriptores adicionales o de representaciones aprendidas explicables sin alterar la interfaz pública del sistema, preservando el mismo régimen de validación y de control de calibración \cite{Amann2019,OConnor2020,BuitragoDuque2023,Pedregosa2011}.
s